% GENERATED FILE -- do not edit by hand.
% Regenerate with: make index   (tools/gen_question_index.py)
\chapter{Question Index}
\label{chap:question_index}

\begin{tldr}
Every one of the 88 interview questions in Volume IV, indexed by chapter, by round, and by frequency. This appendix is generated from the questions themselves, so it cannot drift from the chapters. Work the \freqhigh{} list first: those are the questions that come up in almost every loop.
\end{tldr}

\section{At a Glance}
\begin{table}[H]\centering
\caption{Question distribution}
\begin{tabular}{lr}
\toprule
\textbf{Category} & \textbf{Count} \\
\midrule
\quad \textbf{Total} & \textbf{88} \\
\quad L5 questions & 38 \\
\quad L6 questions & 43 \\
\quad L7 questions & 7 \\
\quad Breadth round & 27 \\
\quad Depth round & 42 \\
\quad Coding round & 10 \\
\quad Design round & 9 \\
\quad \freqhigh{} asked constantly & 35 \\
\quad \freqmed{} common & 45 \\
\quad \freqlow{} occasional & 8 \\
\bottomrule
\end{tabular}
\end{table}

\section{Start Here: The \freqhigh{} Questions}
These come up in almost every loop at this level. If you have one evening, these are the evening.
\begin{enumerate}[leftmargin=*]
\item \textbf{Ch.~1} \quad L5 \quad \textit{Mathematical} \quad Why does L1 regularization produce sparse solutions and L2 does not?
\item \textbf{Ch.~1} \quad L5 \quad \textit{Mathematical} \quad Derive the normal equations for linear regression. When would you not use the closed form?
\item \textbf{Ch.~1} \quad L5 \quad \textit{First Principles} \quad Derive logistic regression's loss from maximum likelihood, give its gradient, and explain why we do not just use squared error
\item \textbf{Ch.~1} \quad L6 \quad \textit{Trade-off} \quad When would you deploy logistic regression instead of gradient boosting in 2026?
\item \textbf{Ch.~2} \quad L5 \quad \textit{Mathematical} \quad Derive the optimal leaf weight and the split-gain formula in XGBoost
\item \textbf{Ch.~2} \quad L5 \quad \textit{Conceptual} \quad Random forest vs.\ gradient boosting---how do they differ, and when do you reach for each?
\item \textbf{Ch.~2} \quad L5 \quad \textit{First Principles} \quad Explain gradient boosting to someone who already understands gradient descent
\item \textbf{Ch.~2} \quad L6 \quad \textit{Trade-off} \quad XGBoost vs.\ LightGBM vs.\ CatBoost: 600K rows, 45 features of which 14 are categorical including a 40K-cardinality \texttt{merchant\_id}. Choose and defend
\item \textbf{Ch.~2} \quad L6 \quad \textit{Conceptual} \quad Why does a gradient-boosted tree beat your MLP on this tabular dataset---and when would it not?
\item \textbf{Ch.~3} \quad L5 \quad \textit{Mathematical} \quad Derive PCA from first principles. Whiteboard it
\item \textbf{Ch.~3} \quad L5 \quad \textit{First Principles} \quad Why does k-means converge, and what does it converge to? And why do you need k-means++?
\item \textbf{Ch.~3} \quad L6 \quad \textit{Mathematical} \quad Derive the EM updates for a Gaussian mixture model, and explain why the likelihood increases monotonically
\item \textbf{Ch.~3} \quad L5 \quad \textit{Trade-off} \quad How do you choose $k$? Your teammate says the elbow plot shows 8 clusters
\item \textbf{Ch.~3} \quad L6 \quad \textit{Debugging} \quad Your t-SNE plot of user embeddings shows five clean, well-separated clusters. What can you conclude?
\item \textbf{Ch.~3} \quad L5 \quad \textit{Conceptual} \quad When does k-means fail, and what do you use instead?
\item \textbf{Ch.~4} \quad L5 \quad \textit{First Principles} \quad MLE, MAP, and fully Bayesian---define each, and tell me where regularization comes from
\item \textbf{Ch.~4} \quad L5 \quad \textit{Debugging} \quad Training RMSE 0.31, validation RMSE 0.94, and the two learning curves are still far apart at 200K rows. Diagnose it and name the knobs, on the model of your choice
\item \textbf{Ch.~4} \quad L6 \quad \textit{Debugging} \quad Our fraud model has 0.92 AUC, but the risk team says the probabilities are useless---their auto-decline rule fires at 0.9 and almost nothing crosses it. What is wrong, and how do you fix it?
\item \textbf{Ch.~5} \quad L6 \quad \textit{Debugging} \quad Your churn model has 0.99 AUC offline and 0.62 in production. Walk me through it
\item \textbf{Ch.~5} \quad L6 \quad \textit{Mathematical} \quad You want to target-encode a 50K-cardinality \texttt{merchant\_id}. Give me the exact construction, and tell me where it leaks if you get it wrong
\item \textbf{Ch.~5} \quad L5 \quad \textit{Debugging} \quad Here is a pipeline. \texttt{StandardScaler} on the full dataframe, then median imputation, then \texttt{SelectKBest(k=100)} on all rows, then SMOTE to balance, then a 5-fold CV that reports 0.94 AUC. Find the bugs and rank them by damage
\item \textbf{Ch.~5} \quad L5 \quad \textit{Conceptual} \quad When is random $k$-fold cross-validation invalid, and what do you use instead?
\item \textbf{Ch.~5} \quad L6 \quad \textit{Trade-off} \quad Fraud is 0.1\% of transactions. Walk me through building the classifier
\item \textbf{Ch.~6} \quad L5 \quad \textit{Estimation} \quad We want to detect a 1\% relative lift on a 2\% click-through rate. How many users do we need, and how long should we run?
\item \textbf{Ch.~6} \quad L6 \quad \textit{Trade-off} \quad Your test came back flat, but you are confident the feature helps. What do you do?
\item \textbf{Ch.~6} \quad L6 \quad \textit{Debugging} \quad Your treatment shows +2\% on the primary metric, but the SRM check fired. What now?
\item \textbf{Ch.~6} \quad L6 \quad \textit{Conceptual} \quad Your marketplace A/B test shows +2\% GMV in treatment. Why might the true effect be smaller, or negative?
\item \textbf{Ch.~7} \quad L6 \quad \textit{System Design} \quad Forecast daily demand for 5{,}000 SKUs across 20 warehouses, 28 days ahead. Walk me through your approach
\item \textbf{Ch.~7} \quad L5 \quad \textit{Conceptual} \quad Why can't you use standard $k$-fold cross-validation on time series, exactly---and what do you use instead?
\item \textbf{Ch.~7} \quad L5 \quad \textit{Debugging} \quad Your GBM forecast flatlines just under the recent maximum while demand keeps growing. What is happening and how do you fix it?
\item \textbf{Ch.~7} \quad L6 \quad \textit{Debugging} \quad Your backtest MASE is 0.72 and everyone is thrilled. In production the forecasts are terrible. Give me your differential diagnosis
\item \textbf{Ch.~8} \quad L5 \quad \textit{Coding} \quad Implement k-means with k-means++ initialization. Handle empty clusters and give me a real convergence criterion
\item \textbf{Ch.~8} \quad L5 \quad \textit{Coding} \quad Implement logistic regression trained by minibatch SGD, derive the gradient, and convince me the gradient is right
\item \textbf{Ch.~8} \quad L6 \quad \textit{Coding} \quad Given a node's samples and labels, find the split that maximizes impurity decrease. What is the complexity of your approach?
\item \textbf{Ch.~8} \quad L5 \quad \textit{Coding} \quad Implement k-nearest-neighbors classification. Make the query as fast as you can without a tree
\end{enumerate}

\section{By Chapter}
\subsection*{Chapter 1: Supervised Learning Core}
\begin{tabularx}{\textwidth}{@{}l l l X@{}}
\toprule
\textbf{Lvl} & \textbf{Freq} & \textbf{Type} & \textbf{Question} \\
\midrule
L5 & \freqhigh{} & Mathematical & Why does L1 regularization produce sparse solutions and L2 does not? \\
L5 & \freqhigh{} & Mathematical & Derive the normal equations for linear regression. When would you not use the closed form? \\
L5 & \freqhigh{} & First Principles & Derive logistic regression's loss from maximum likelihood, give its gradient, and explain why we do not just use squared error \\
L6 & \freqmed{} & Mathematical & Derive the SVM: margin, primal, dual, and the kernel trick. What exactly is a support vector? \\
L6 & \freqmed{} & Conceptual & Generative versus discriminative---define them, and use Naive Bayes and logistic regression to explain the trade-off \\
L5 & \freqmed{} & Debugging & Your logistic regression's weights blew up to $\pm 10^6$ and training accuracy is exactly 1.0. What happened, and what do you do? \\
L5 & \freqmed{} & Conceptual & What does $C$ do in an SVM? What about $\gamma$ in an RBF kernel? \\
L5 & \freqmed{} & Conceptual & Why does $k$-NN degrade in high dimensions? Quantify it \\
L6 & \freqhigh{} & Trade-off & When would you deploy logistic regression instead of gradient boosting in 2026? \\
L6 & \freqmed{} & Debugging & You fit a linear model on 150 rows and 200 features. The coefficients flip sign between refits and two features have VIF over 40. Walk me through it \\
L6 & \freqlow{} & Architecture Design & You need to model weekly support tickets per account. The target is a non-negative integer, most accounts have zero or one, a few have hundreds, and accounts have very different sizes. Design the model \\
L7 & \freqlow{} & System Design & Design the model layer for an ad click-through-rate system: $10^9$ events per day, $10^8$ hashed sparse features, hourly retraining, sub-millisecond scoring, and the score is multiplied by a bid \\
\bottomrule
\end{tabularx}

\subsection*{Chapter 2: Trees and Ensembles}
\begin{tabularx}{\textwidth}{@{}l l l X@{}}
\toprule
\textbf{Lvl} & \textbf{Freq} & \textbf{Type} & \textbf{Question} \\
\midrule
L5 & \freqhigh{} & Mathematical & Derive the optimal leaf weight and the split-gain formula in XGBoost \\
L5 & \freqhigh{} & Conceptual & Random forest vs.\ gradient boosting---how do they differ, and when do you reach for each? \\
L5 & \freqhigh{} & First Principles & Explain gradient boosting to someone who already understands gradient descent \\
L6 & \freqhigh{} & Trade-off & XGBoost vs.\ LightGBM vs.\ CatBoost: 600K rows, 45 features of which 14 are categorical including a 40K-cardinality \texttt{merchant\_id}. Choose and defend \\
L6 & \freqhigh{} & Conceptual & Why does a gradient-boosted tree beat your MLP on this tabular dataset---and when would it not? \\
L6 & \freqmed{} & Debugging & Your LambdaMART ranker's training NDCG keeps improving, but validation NDCG peaked at tree 200 and is now degrading. Walk me through your response \\
L5 & \freqmed{} & Conceptual & Why do trees split on Gini or entropy instead of misclassification error, since error is what we ultimately care about? \\
L5 & \freqmed{} & Mathematical & What fraction of the training data does each tree in a random forest see? Derive it, and explain what the rest is used for \\
L6 & \freqmed{} & Conceptual & How does CatBoost avoid target leakage with categorical features, and what is ``ordered boosting'' fixing? \\
L5 & \freqmed{} & Conceptual & Why can't a tree-based model extrapolate, and where does that bite in practice? \\
L6 & \freqmed{} & Debugging & Your random forest's impurity-based feature importance says \texttt{session\_id\_hash} is the top feature. The PM wants to build product strategy around the ranking. What do you tell them? \\
L7 & \freqlow{} & System Design & Design the serving story for a GBDT fraud model: p99 model latency under 2\,ms at 20K QPS, with a compliance requirement that risk never decreases as chargeback count increases \\
\bottomrule
\end{tabularx}

\subsection*{Chapter 3: Unsupervised Learning}
\begin{tabularx}{\textwidth}{@{}l l l X@{}}
\toprule
\textbf{Lvl} & \textbf{Freq} & \textbf{Type} & \textbf{Question} \\
\midrule
L5 & \freqhigh{} & Mathematical & Derive PCA from first principles. Whiteboard it \\
L5 & \freqhigh{} & First Principles & Why does k-means converge, and what does it converge to? And why do you need k-means++? \\
L6 & \freqhigh{} & Mathematical & Derive the EM updates for a Gaussian mixture model, and explain why the likelihood increases monotonically \\
L5 & \freqhigh{} & Trade-off & How do you choose $k$? Your teammate says the elbow plot shows 8 clusters \\
L6 & \freqhigh{} & Debugging & Your t-SNE plot of user embeddings shows five clean, well-separated clusters. What can you conclude? \\
L5 & \freqhigh{} & Conceptual & When does k-means fail, and what do you use instead? \\
L5 & \freqmed{} & Mathematical & What is the relationship between PCA and SVD, and why do libraries use the SVD? \\
L6 & \freqmed{} & Mathematical & How is k-means a special case of a Gaussian mixture model? \\
L6 & \freqmed{} & Trade-off & You have 10K-dimensional sparse features and want to reduce them for a downstream classifier. PCA, or something else? \\
L6 & \freqlow{} & Debugging & Your GMM's training log-likelihood spikes toward infinity and one component's weight collapses. What happened and what do you do? \\
L7 & \freqlow{} & System Design & Design a weekly user-segmentation system: 80M users, roughly 200 behavioural features, output consumed by marketing and by a personalization service \\
\bottomrule
\end{tabularx}

\subsection*{Chapter 4: Probabilistic Foundations}
\begin{tabularx}{\textwidth}{@{}l l l X@{}}
\toprule
\textbf{Lvl} & \textbf{Freq} & \textbf{Type} & \textbf{Question} \\
\midrule
L5 & \freqhigh{} & First Principles & MLE, MAP, and fully Bayesian---define each, and tell me where regularization comes from \\
L5 & \freqmed{} & Mathematical & Show that ridge regression is MAP estimation. What is $\lambda$ in terms of the model? \\
L5 & \freqmed{} & Mathematical & Why is the MLE of the Gaussian variance biased? Does it matter? \\
L5 & \freqhigh{} & Debugging & Training RMSE 0.31, validation RMSE 0.94, and the two learning curves are still far apart at 200K rows. Diagnose it and name the knobs, on the model of your choice \\
L6 & \freqhigh{} & Debugging & Our fraud model has 0.92 AUC, but the risk team says the probabilities are useless---their auto-decline rule fires at 0.9 and almost nothing crosses it. What is wrong, and how do you fix it? \\
L6 & \freqmed{} & Mathematical & Prove that log loss is a proper scoring rule. Why should I care? \\
L5 & \freqmed{} & Conceptual & Which classical models give you calibrated probabilities out of the box, and which do not? \\
L5 & \freqmed{} & Conceptual & Confidence interval versus prediction interval---define both, and tell me which one the product manager actually wanted \\
L6 & \freqmed{} & First Principles & I want prediction sets with a coverage guarantee and no distributional assumptions. Build it, then tell me when the guarantee is worthless \\
L6 & \freqlow{} & Trade-off & Logistics wants a P90 delivery-time estimate per order, not a point prediction. What do you build, and why is your loss function correct? \\
L7 & \freqmed{} & System Design & We are launching an automated credit-limit model. Design the uncertainty layer: what we estimate, how, and what we monitor \\
\bottomrule
\end{tabularx}

\subsection*{Chapter 5: The Applied Craft}
\begin{tabularx}{\textwidth}{@{}l l l X@{}}
\toprule
\textbf{Lvl} & \textbf{Freq} & \textbf{Type} & \textbf{Question} \\
\midrule
L6 & \freqhigh{} & Debugging & Your churn model has 0.99 AUC offline and 0.62 in production. Walk me through it \\
L6 & \freqhigh{} & Mathematical & You want to target-encode a 50K-cardinality \texttt{merchant\_id}. Give me the exact construction, and tell me where it leaks if you get it wrong \\
L5 & \freqhigh{} & Debugging & Here is a pipeline. \texttt{StandardScaler} on the full dataframe, then median imputation, then \texttt{SelectKBest(k=100)} on all rows, then SMOTE to balance, then a 5-fold CV that reports 0.94 AUC. Find the bugs and rank them by damage \\
L5 & \freqhigh{} & Conceptual & When is random $k$-fold cross-validation invalid, and what do you use instead? \\
L6 & \freqhigh{} & Trade-off & Fraud is 0.1\% of transactions. Walk me through building the classifier \\
L6 & \freqmed{} & Debugging & Your teammate added SMOTE and offline $F_1$ improved, but the production fraud model got worse. Explain \\
L6 & \freqmed{} & Conceptual & What is nested cross-validation, and when do you actually need it? \\
L6 & \freqmed{} & Conceptual & Two features are highly correlated. What happens to permutation importance and to SHAP? \\
L6 & \freqmed{} & Trade-off & A PM sees that \texttt{discount\_rate} has the highest SHAP value in the conversion model and asks whether we should raise discounts. What do you say? \\
L5 & \freqmed{} & Conceptual & A key feature is missing for 30\% of rows. What do you do? \\
L6 & \freqmed{} & Estimation & 100M transactions a day, 0.1\% fraud, and a review team that can handle 500 alerts a day. What precision and recall can you promise? \\
L7 & \freqlow{} & System Design & Design the offline evaluation protocol for a churn model that will be retrained weekly and must survive an audit \\
\bottomrule
\end{tabularx}

\subsection*{Chapter 6: Experimentation and Causal Inference}
\begin{tabularx}{\textwidth}{@{}l l l X@{}}
\toprule
\textbf{Lvl} & \textbf{Freq} & \textbf{Type} & \textbf{Question} \\
\midrule
L5 & \freqhigh{} & Estimation & We want to detect a 1\% relative lift on a 2\% click-through rate. How many users do we need, and how long should we run? \\
L6 & \freqhigh{} & Trade-off & Your test came back flat, but you are confident the feature helps. What do you do? \\
L6 & \freqhigh{} & Debugging & Your treatment shows +2\% on the primary metric, but the SRM check fired. What now? \\
L6 & \freqmed{} & Conceptual & Explain CUPED to a PM, and quantify what it buys us \\
L6 & \freqhigh{} & Conceptual & Your marketplace A/B test shows +2\% GMV in treatment. Why might the true effect be smaller, or negative? \\
L6 & \freqmed{} & Conceptual & Why not send retention offers to the users with the highest churn probability? \\
L6 & \freqmed{} & Trade-off & Thompson sampling, UCB, or $\epsilon$-greedy---mechanics, and when would you use a bandit instead of an A/B test? \\
L7 & \freqmed{} & Architecture Design & We cannot randomize prices. How would you estimate price elasticity from historical data? \\
L5 & \freqmed{} & Conceptual & What assumption does difference-in-differences require, and how would you probe it? \\
L7 & \freqmed{} & Architecture Design & Design the experimentation platform for a 200-engineer organization \\
\bottomrule
\end{tabularx}

\subsection*{Chapter 7: Time Series Essentials}
\begin{tabularx}{\textwidth}{@{}l l l X@{}}
\toprule
\textbf{Lvl} & \textbf{Freq} & \textbf{Type} & \textbf{Question} \\
\midrule
L6 & \freqhigh{} & System Design & Forecast daily demand for 5{,}000 SKUs across 20 warehouses, 28 days ahead. Walk me through your approach \\
L5 & \freqhigh{} & Conceptual & Why can't you use standard $k$-fold cross-validation on time series, exactly---and what do you use instead? \\
L5 & \freqhigh{} & Debugging & Your GBM forecast flatlines just under the recent maximum while demand keeps growing. What is happening and how do you fix it? \\
L6 & \freqhigh{} & Debugging & Your backtest MASE is 0.72 and everyone is thrilled. In production the forecasts are terrible. Give me your differential diagnosis \\
L5 & \freqmed{} & Mathematical & You need one accuracy number across 5{,}000 SKUs of wildly different volume. Which metric, and why not MAPE? \\
L5 & \freqmed{} & Conceptual & What is stationarity, and why do ARIMA-type models want it? \\
L6 & \freqmed{} & Trade-off & Recursive or direct multi-step forecasting---which do you pick and why? \\
L6 & \freqmed{} & First Principles & Build me a forecast with uncertainty bands the business will actually trust \\
L6 & \freqmed{} & Trade-off & When would you reach for a deep sequence model instead of ETS or a global GBM? \\
L5 & \freqlow{} & Mathematical & I show you an ACF that decays geometrically and a PACF with one big spike at lag 1 and nothing after. What model, and what does the parameter mean? \\
\bottomrule
\end{tabularx}

\subsection*{Chapter 8: The From-Scratch Coding Canon}
\begin{tabularx}{\textwidth}{@{}l l l X@{}}
\toprule
\textbf{Lvl} & \textbf{Freq} & \textbf{Type} & \textbf{Question} \\
\midrule
L5 & \freqhigh{} & Coding & Implement k-means with k-means++ initialization. Handle empty clusters and give me a real convergence criterion \\
L5 & \freqhigh{} & Coding & Implement logistic regression trained by minibatch SGD, derive the gradient, and convince me the gradient is right \\
L6 & \freqhigh{} & Coding & Given a node's samples and labels, find the split that maximizes impurity decrease. What is the complexity of your approach? \\
L6 & \freqmed{} & Coding & Implement gradient boosting for binary classification with stumps. Then tell me where the leaf values come from \\
L6 & \freqmed{} & Coding & Implement PCA. Use power iteration rather than calling an eigensolver, and return the top $k$ components \\
L5 & \freqhigh{} & Coding & Implement k-nearest-neighbors classification. Make the query as fast as you can without a tree \\
L6 & \freqmed{} & Coding & Now build a kd-tree and use it for nearest-neighbor search. When does it stop helping? \\
L5 & \freqmed{} & Coding & Implement linear regression two ways: the closed form and gradient descent. Which would a library use? \\
L5 & \freqmed{} & Coding & Implement AUC from scratch. What do you do about ties? \\
L6 & \freqmed{} & Coding & Write a cross-validation harness without sklearn, and use it to compare two models. Where does leakage sneak in? \\
\bottomrule
\end{tabularx}

\section{By Round}
The rounds a loop is actually made of. Drill one column at a time.
\subsection*{Breadth (27 questions)}
\begin{itemize}[leftmargin=*]
\item \textbf{Ch.~1} \quad L5 \quad Derive logistic regression's loss from maximum likelihood, give its gradient, and explain why we do not just use squared error
\item \textbf{Ch.~1} \quad L6 \quad Generative versus discriminative---define them, and use Naive Bayes and logistic regression to explain the trade-off
\item \textbf{Ch.~1} \quad L5 \quad What does $C$ do in an SVM? What about $\gamma$ in an RBF kernel?
\item \textbf{Ch.~1} \quad L5 \quad Why does $k$-NN degrade in high dimensions? Quantify it
\item \textbf{Ch.~2} \quad L5 \quad Random forest vs.\ gradient boosting---how do they differ, and when do you reach for each?
\item \textbf{Ch.~2} \quad L5 \quad Explain gradient boosting to someone who already understands gradient descent
\item \textbf{Ch.~2} \quad L6 \quad Why does a gradient-boosted tree beat your MLP on this tabular dataset---and when would it not?
\item \textbf{Ch.~2} \quad L5 \quad Why do trees split on Gini or entropy instead of misclassification error, since error is what we ultimately care about?
\item \textbf{Ch.~2} \quad L6 \quad How does CatBoost avoid target leakage with categorical features, and what is ``ordered boosting'' fixing?
\item \textbf{Ch.~2} \quad L5 \quad Why can't a tree-based model extrapolate, and where does that bite in practice?
\item \textbf{Ch.~3} \quad L5 \quad Why does k-means converge, and what does it converge to? And why do you need k-means++?
\item \textbf{Ch.~3} \quad L5 \quad When does k-means fail, and what do you use instead?
\item \textbf{Ch.~4} \quad L5 \quad MLE, MAP, and fully Bayesian---define each, and tell me where regularization comes from
\item \textbf{Ch.~4} \quad L5 \quad Which classical models give you calibrated probabilities out of the box, and which do not?
\item \textbf{Ch.~4} \quad L5 \quad Confidence interval versus prediction interval---define both, and tell me which one the product manager actually wanted
\item \textbf{Ch.~4} \quad L6 \quad I want prediction sets with a coverage guarantee and no distributional assumptions. Build it, then tell me when the guarantee is worthless
\item \textbf{Ch.~5} \quad L5 \quad When is random $k$-fold cross-validation invalid, and what do you use instead?
\item \textbf{Ch.~5} \quad L6 \quad What is nested cross-validation, and when do you actually need it?
\item \textbf{Ch.~5} \quad L6 \quad Two features are highly correlated. What happens to permutation importance and to SHAP?
\item \textbf{Ch.~5} \quad L5 \quad A key feature is missing for 30\% of rows. What do you do?
\item \textbf{Ch.~6} \quad L6 \quad Explain CUPED to a PM, and quantify what it buys us
\item \textbf{Ch.~6} \quad L6 \quad Your marketplace A/B test shows +2\% GMV in treatment. Why might the true effect be smaller, or negative?
\item \textbf{Ch.~6} \quad L6 \quad Why not send retention offers to the users with the highest churn probability?
\item \textbf{Ch.~6} \quad L5 \quad What assumption does difference-in-differences require, and how would you probe it?
\item \textbf{Ch.~7} \quad L5 \quad Why can't you use standard $k$-fold cross-validation on time series, exactly---and what do you use instead?
\item \textbf{Ch.~7} \quad L5 \quad What is stationarity, and why do ARIMA-type models want it?
\item \textbf{Ch.~7} \quad L6 \quad Build me a forecast with uncertainty bands the business will actually trust
\end{itemize}

\subsection*{Depth (42 questions)}
\begin{itemize}[leftmargin=*]
\item \textbf{Ch.~1} \quad L5 \quad Why does L1 regularization produce sparse solutions and L2 does not?
\item \textbf{Ch.~1} \quad L5 \quad Derive the normal equations for linear regression. When would you not use the closed form?
\item \textbf{Ch.~1} \quad L6 \quad Derive the SVM: margin, primal, dual, and the kernel trick. What exactly is a support vector?
\item \textbf{Ch.~1} \quad L5 \quad Your logistic regression's weights blew up to $\pm 10^6$ and training accuracy is exactly 1.0. What happened, and what do you do?
\item \textbf{Ch.~1} \quad L6 \quad When would you deploy logistic regression instead of gradient boosting in 2026?
\item \textbf{Ch.~1} \quad L6 \quad You fit a linear model on 150 rows and 200 features. The coefficients flip sign between refits and two features have VIF over 40. Walk me through it
\item \textbf{Ch.~2} \quad L5 \quad Derive the optimal leaf weight and the split-gain formula in XGBoost
\item \textbf{Ch.~2} \quad L6 \quad XGBoost vs.\ LightGBM vs.\ CatBoost: 600K rows, 45 features of which 14 are categorical including a 40K-cardinality \texttt{merchant\_id}. Choose and defend
\item \textbf{Ch.~2} \quad L6 \quad Your LambdaMART ranker's training NDCG keeps improving, but validation NDCG peaked at tree 200 and is now degrading. Walk me through your response
\item \textbf{Ch.~2} \quad L5 \quad What fraction of the training data does each tree in a random forest see? Derive it, and explain what the rest is used for
\item \textbf{Ch.~2} \quad L6 \quad Your random forest's impurity-based feature importance says \texttt{session\_id\_hash} is the top feature. The PM wants to build product strategy around the ranking. What do you tell them?
\item \textbf{Ch.~3} \quad L5 \quad Derive PCA from first principles. Whiteboard it
\item \textbf{Ch.~3} \quad L6 \quad Derive the EM updates for a Gaussian mixture model, and explain why the likelihood increases monotonically
\item \textbf{Ch.~3} \quad L5 \quad How do you choose $k$? Your teammate says the elbow plot shows 8 clusters
\item \textbf{Ch.~3} \quad L6 \quad Your t-SNE plot of user embeddings shows five clean, well-separated clusters. What can you conclude?
\item \textbf{Ch.~3} \quad L5 \quad What is the relationship between PCA and SVD, and why do libraries use the SVD?
\item \textbf{Ch.~3} \quad L6 \quad How is k-means a special case of a Gaussian mixture model?
\item \textbf{Ch.~3} \quad L6 \quad You have 10K-dimensional sparse features and want to reduce them for a downstream classifier. PCA, or something else?
\item \textbf{Ch.~3} \quad L6 \quad Your GMM's training log-likelihood spikes toward infinity and one component's weight collapses. What happened and what do you do?
\item \textbf{Ch.~4} \quad L5 \quad Show that ridge regression is MAP estimation. What is $\lambda$ in terms of the model?
\item \textbf{Ch.~4} \quad L5 \quad Why is the MLE of the Gaussian variance biased? Does it matter?
\item \textbf{Ch.~4} \quad L5 \quad Training RMSE 0.31, validation RMSE 0.94, and the two learning curves are still far apart at 200K rows. Diagnose it and name the knobs, on the model of your choice
\item \textbf{Ch.~4} \quad L6 \quad Our fraud model has 0.92 AUC, but the risk team says the probabilities are useless---their auto-decline rule fires at 0.9 and almost nothing crosses it. What is wrong, and how do you fix it?
\item \textbf{Ch.~4} \quad L6 \quad Prove that log loss is a proper scoring rule. Why should I care?
\item \textbf{Ch.~4} \quad L6 \quad Logistics wants a P90 delivery-time estimate per order, not a point prediction. What do you build, and why is your loss function correct?
\item \textbf{Ch.~5} \quad L6 \quad Your churn model has 0.99 AUC offline and 0.62 in production. Walk me through it
\item \textbf{Ch.~5} \quad L6 \quad You want to target-encode a 50K-cardinality \texttt{merchant\_id}. Give me the exact construction, and tell me where it leaks if you get it wrong
\item \textbf{Ch.~5} \quad L5 \quad Here is a pipeline. \texttt{StandardScaler} on the full dataframe, then median imputation, then \texttt{SelectKBest(k=100)} on all rows, then SMOTE to balance, then a 5-fold CV that reports 0.94 AUC. Find the bugs and rank them by damage
\item \textbf{Ch.~5} \quad L6 \quad Fraud is 0.1\% of transactions. Walk me through building the classifier
\item \textbf{Ch.~5} \quad L6 \quad Your teammate added SMOTE and offline $F_1$ improved, but the production fraud model got worse. Explain
\item \textbf{Ch.~5} \quad L6 \quad A PM sees that \texttt{discount\_rate} has the highest SHAP value in the conversion model and asks whether we should raise discounts. What do you say?
\item \textbf{Ch.~5} \quad L6 \quad 100M transactions a day, 0.1\% fraud, and a review team that can handle 500 alerts a day. What precision and recall can you promise?
\item \textbf{Ch.~6} \quad L5 \quad We want to detect a 1\% relative lift on a 2\% click-through rate. How many users do we need, and how long should we run?
\item \textbf{Ch.~6} \quad L6 \quad Your test came back flat, but you are confident the feature helps. What do you do?
\item \textbf{Ch.~6} \quad L6 \quad Your treatment shows +2\% on the primary metric, but the SRM check fired. What now?
\item \textbf{Ch.~6} \quad L6 \quad Thompson sampling, UCB, or $\epsilon$-greedy---mechanics, and when would you use a bandit instead of an A/B test?
\item \textbf{Ch.~7} \quad L5 \quad Your GBM forecast flatlines just under the recent maximum while demand keeps growing. What is happening and how do you fix it?
\item \textbf{Ch.~7} \quad L6 \quad Your backtest MASE is 0.72 and everyone is thrilled. In production the forecasts are terrible. Give me your differential diagnosis
\item \textbf{Ch.~7} \quad L5 \quad You need one accuracy number across 5{,}000 SKUs of wildly different volume. Which metric, and why not MAPE?
\item \textbf{Ch.~7} \quad L6 \quad Recursive or direct multi-step forecasting---which do you pick and why?
\item \textbf{Ch.~7} \quad L6 \quad When would you reach for a deep sequence model instead of ETS or a global GBM?
\item \textbf{Ch.~7} \quad L5 \quad I show you an ACF that decays geometrically and a PACF with one big spike at lag 1 and nothing after. What model, and what does the parameter mean?
\end{itemize}

\subsection*{Coding (10 questions)}
\begin{itemize}[leftmargin=*]
\item \textbf{Ch.~8} \quad L5 \quad Implement k-means with k-means++ initialization. Handle empty clusters and give me a real convergence criterion
\item \textbf{Ch.~8} \quad L5 \quad Implement logistic regression trained by minibatch SGD, derive the gradient, and convince me the gradient is right
\item \textbf{Ch.~8} \quad L6 \quad Given a node's samples and labels, find the split that maximizes impurity decrease. What is the complexity of your approach?
\item \textbf{Ch.~8} \quad L6 \quad Implement gradient boosting for binary classification with stumps. Then tell me where the leaf values come from
\item \textbf{Ch.~8} \quad L6 \quad Implement PCA. Use power iteration rather than calling an eigensolver, and return the top $k$ components
\item \textbf{Ch.~8} \quad L5 \quad Implement k-nearest-neighbors classification. Make the query as fast as you can without a tree
\item \textbf{Ch.~8} \quad L6 \quad Now build a kd-tree and use it for nearest-neighbor search. When does it stop helping?
\item \textbf{Ch.~8} \quad L5 \quad Implement linear regression two ways: the closed form and gradient descent. Which would a library use?
\item \textbf{Ch.~8} \quad L5 \quad Implement AUC from scratch. What do you do about ties?
\item \textbf{Ch.~8} \quad L6 \quad Write a cross-validation harness without sklearn, and use it to compare two models. Where does leakage sneak in?
\end{itemize}

\subsection*{Design (9 questions)}
\begin{itemize}[leftmargin=*]
\item \textbf{Ch.~1} \quad L6 \quad You need to model weekly support tickets per account. The target is a non-negative integer, most accounts have zero or one, a few have hundreds, and accounts have very different sizes. Design the model
\item \textbf{Ch.~1} \quad L7 \quad Design the model layer for an ad click-through-rate system: $10^9$ events per day, $10^8$ hashed sparse features, hourly retraining, sub-millisecond scoring, and the score is multiplied by a bid
\item \textbf{Ch.~2} \quad L7 \quad Design the serving story for a GBDT fraud model: p99 model latency under 2\,ms at 20K QPS, with a compliance requirement that risk never decreases as chargeback count increases
\item \textbf{Ch.~3} \quad L7 \quad Design a weekly user-segmentation system: 80M users, roughly 200 behavioural features, output consumed by marketing and by a personalization service
\item \textbf{Ch.~4} \quad L7 \quad We are launching an automated credit-limit model. Design the uncertainty layer: what we estimate, how, and what we monitor
\item \textbf{Ch.~5} \quad L7 \quad Design the offline evaluation protocol for a churn model that will be retrained weekly and must survive an audit
\item \textbf{Ch.~6} \quad L7 \quad We cannot randomize prices. How would you estimate price elasticity from historical data?
\item \textbf{Ch.~6} \quad L7 \quad Design the experimentation platform for a 200-engineer organization
\item \textbf{Ch.~7} \quad L6 \quad Forecast daily demand for 5{,}000 SKUs across 20 warehouses, 28 days ahead. Walk me through your approach
\end{itemize}

